% PAGES: 74-76
% Continues section 2.5.1 "LU linear solvers for tridiagonal matrices"
% (subsection of 2.5 "Linear solvers using LU decomposition without
% pivoting"), begun before page 74 in a sibling agent's file. Do not open
% that file. This file ends mid-sentence at the bottom of page 76; the
% sentence is completed at the top of page 77 in sec02_c4_2.tex.

Moreover, since
\[
U(i,i+1:n) \;=\; \bigl(U(i,i+1)\ \ 0\ \ \ldots\ \ 0\bigr)
\]
\[
\text{and} \quad L(i+1:n,i) \;=\; \begin{pmatrix} L(i+1,i)\\ 0\\ \vdots\\ 0\end{pmatrix},
\]
we find that
\begin{align*}
L(i+1:n,i)\,U(i,i+1:n)
&\;=\; \begin{pmatrix} L(i+1,i)\\ 0\\ \vdots\\ 0\end{pmatrix}
\bigl(U(i,i+1)\ \ 0\ \ \ldots\ \ 0\bigr)\\[6pt]
&\;=\; \begin{pmatrix}
L(i+1,i)U(i,i+1) & 0 & \ldots & 0\\
0 & 0 & \ldots & 0\\
\vdots & \vdots & \ddots & \vdots\\
0 & 0 & \ldots & 0
\end{pmatrix}.
\end{align*}

Then, updating the $(n-i)\times(n-i)$ lower right part of $A$ as in (2.37), i.e.,
\[
A(i+1:n,i+1:n) \;=\; A(i+1:n,i+1:n) \;-\; L(i+1:n,i)\,U(i,i+1:n),
\]
involves only changing the value of the $A(i+1,i+1)$ entry to
\[
A(i+1,i+1) \;=\; A(i+1,i+1) \;-\; L(i+1,i)U(i,i+1),
\]
which preserves the tridiagonal structure of the updated $(n-i-1)\times(n-i-1)$
matrix $A(i+1:n,i+1:n)$.

The LU decomposition routine from Table 2.5 simplifies to the routine from Table 2.7
for tridiagonal matrices with LU decomposition without pivoting.
The operation count for the tridiagonal LU decomposition without pivoting from
Table 2.7 is $3n-3$ operations: in each step of the ``for'' loop

\begin{center}
\fbox{\begin{minipage}{0.82\textwidth}
\raggedright
for $i = 1 : (n-1)$\\
\hspace*{1.5em}$L(i,i) = 1;\ L(i+1,i) = A(i+1,i)/A(i,i);$\\
\hspace*{1.5em}$U(i,i) = A(i,i);\ U(i,i+1) = A(i,i+1);$\\
\hspace*{1.5em}$A(i+1,i+1) = A(i+1,i+1) - L(i+1,i)U(i,i+1);$\\
end
\end{minipage}}
\end{center}

we perform $3$ operations, for a total number of operations of
\begin{equation}
3(n-1) \;=\; 3n-3.
\end{equation}

Thus, if $A$ is a tridiagonal matrix with LU decomposition without pivoting, the
LU decomposition $[L,U] = lu\_no\_pivoting(A)$ from the linear solver routine from
Table 2.6 can be replaced by the solver $[L,U] = lu\_no\_pivoting\_tridiag(A)$ from
Table 2.7.
Moreover, since the matrix $L$ is lower triangular bidiagonal, the forward
substitution $y = forward\_subst(L,b)$ from Table 2.6 can be replaced by the
forward substitution $y = forward\_subst\_bidiag(L,b)$ for lower triangular
bidiagonal matrices, see Table 2.2, and, since the matrix $U$ is upper triangular
bidiagonal, the backward substitution $x = backward\_subst(U,y)$ from Table 2.6 can
be replaced by the backward substitution $x = backward\_subst\_bidiag(U,y)$ for
upper triangular bidiagonal matrices, see Table 2.4.
The resulting routine for solving linear systems corresponding to tridiagonal
matrices using LU decomposition without pivoting can be found in Table 2.8.

\begin{table}[H]
\centering
\caption{LU decomposition without pivoting for tridiagonal matrices}
\fbox{\begin{minipage}{0.88\textwidth}
\raggedright
Function Call:\\
$[L,U] = lu\_no\_pivoting\_tridiag(A)$

\medskip
Input:\\
$A = $ nonsingular tridiagonal matrix of size $n$ with LU decomposition

\medskip
Output:\\
$L = $ lower triangular bidiagonal matrix with entries $1$ on main diagonal\\
$U = $ upper triangular bidiagonal matrix\\
such that $A = LU$

\medskip
for $i = 1 : (n-1)$\\
\hspace*{1.5em}$L(i,i) = 1;\ L(i+1,i) = A(i+1,i)/A(i,i);$\\
\hspace*{1.5em}$U(i,i) = A(i,i);\ U(i,i+1) = A(i,i+1);$\\
\hspace*{1.5em}$A(i+1,i+1) = A(i+1,i+1) - L(i+1,i)U(i,i+1);$\\
end\\
$L(n,n) = 1;\ U(n,n) = A(n,n)$
\end{minipage}}
\end{table}

\begin{table}[H]
\centering
\caption{Tridiagonal linear solver using LU decomposition without pivoting}
\fbox{\begin{minipage}{0.88\textwidth}
\raggedright
Function Call:\\
$x = linear\_solve\_LU\_no\_pivoting\_tridiag(A,b)$

\medskip
Input:\\
$A = $ nonsingular tridiagonal matrix of size $n$ with LU decomposition\\
$b = $ column vector of size $n$

\medskip
Output:\\
$x = $ solution to $Ax=b$

\medskip
\small
\begin{tabular}{@{}p{0.60\linewidth}@{\hspace{0.8em}}p{0.35\linewidth}@{}}
$[L,U] = lu\_no\_pivoting\_tridiag(A);$ & \textit{// LU decomposition of $A$}\\
$y = forward\_subst\_bidiag(L,b);$ & \textit{// solve $Ly=b$}\\
$x = backward\_subst\_bidiag(U,y);$ & \textit{// solve $Ux=y$}\\
\end{tabular}
\end{minipage}}
\end{table}

Using the explicit implementations of $lu\_no\_pivoting\_tridiag$ from Table 2.7, of
$forward\_subst\_bidiag$ from Table 2.2, and of $backward\_subst\_bidiag$ from
Table 2.4, the explicit implementation of the tridiagonal linear solver from
Table 2.8 can be found in Table 2.9.
Note that, in the pseudocode from Table 2.9, the following further simplifications

\begin{table}[H]
\centering
\caption{Explicit tridiagonal linear solver using LU decomposition without pivoting}
\fbox{\begin{minipage}{0.88\textwidth}
\raggedright
Function Call:\\
$x = linear\_solve\_LU\_no\_pivoting\_tridiag(A,b)$

\medskip
Input:\\
$A = $ nonsingular tridiagonal matrix of size $n$ with LU decomposition\\
$b = $ column vector of size $n$

\medskip
Output:\\
$x = $ solution to $Ax=b$

\medskip
\small
\begin{tabular}{@{}p{0.58\linewidth}@{\hspace{0.8em}}p{0.37\linewidth}@{}}
for $i = 1 : (n-1)$ & \\
\quad $L(i,i) = 1;\ L(i+1,i) = A(i+1,i)/A(i,i);$ & \\
\quad $U(i,i) = A(i,i);\ U(i,i+1) = A(i,i+1);$ & \\
\quad $A(i+1,i+1) = A(i+1,i+1) - L(i+1,i)U(i,i+1);$ & \\
end & \\
$L(n,n) = 1;\ U(n,n) = A(n,n);$ & \textit{// LU decomposition of $A$}\\
$y(1) = b(1);$ & \\
for $j = 2 : n$ & \\
\quad $y(j) = b(j) - L(j,j-1)y(j-1);$ & \textit{// forward substitution for $Ly=b$}\\
end & \\
$x(n) = \dfrac{y(n)}{U(n,n)};$ & \\
for $j = (n-1) : 1$ & \\
\quad $x(j) = \dfrac{y(j)-U(j,j+1)x(j+1)}{U(j,j)};$ & \\
end & \textit{// backward substitution for $Ux=y$}\\
\end{tabular}
\end{minipage}}
\end{table}

in the explicit implementation of $y = forward\_subst\_bidiag(L,b)$, see Table 2.2,
are included:
\begin{itemize}
\item since $L(1,1)=1$, $x(1) = \dfrac{b(1)}{L(1,1)}$ from Table 2.2 becomes
$x(1) = b(1)$;
\item since $L(j,j)=1$, $x(j) = \dfrac{b(j)-L(j,j-1)x(j-1)}{L(j,j)}$ from Table 2.2
becomes
\[
x(j) \;=\; b(j) - L(j,j-1)x(j-1).
\]
\end{itemize}

Thus, the operation count for the forward substitution ``for'' loop

\begin{center}
\fbox{\begin{minipage}{0.82\textwidth}
\raggedright
for $j = 2 : n$\\
\hspace*{1.5em}$x(j) = b(j) - L(j,j-1)x(j-1);$\\
end
\end{minipage}}
\end{center}

from Table 2.9 is
\begin{equation}
2(n-1) \;=\; 2n-2.
\end{equation}

The operation count for the pseudocode from Table 2.9 is as follows:
\begin{itemize}
\item $3n-3$ for the LU decomposition of $A$; cf. (2.48);
\item $2n-2$ for the forward substitution for $Ly=b$; cf. (2.49);
\item $3n-2$ for the backward substitution for $Ux=y$; cf. (2.20),
\end{itemize}
for a total operation count of
\[
(3n-3) + (2n-2) + (3n-2) \;=\; 8n-7 \;=\; 8n+O(1);
\]
